
\documentclass[9pt]{beamer}
\usepackage{hyperref}
\usepackage{latexsym,xcolor,multicol,booktabs}
\usepackage{amsmath,amssymb}		
\usepackage{graphicx}      
\usepackage{subcaption}
\usepackage{comment}





\title{Homework 1: Robust Root-Finding Solvers}
\author{Yuxuan Zhao}
\date{\today}
\institute{Department of Economics, University of Minnesota}

\usepackage{PKU}


\def\cmd#1{\texttt{\color{red}\footnotesize $\backslash$#1}}
\def\env#1{\texttt{\color{blue}\footnotesize #1}}


\newtheorem{thm}{Theorem}[theorem]


%—-------------------------------------------------------------
% 正文区

\begin{document}
	
	\begin{frame}
		\titlepage
		%\begin{figure}[htpb]
		%	\begin{center}
		%		\includegraphics[width=0.2\linewidth]{PKU-logo.png} 
		%	\end{center}
		%\end{figure}
	\end{frame}
	
	% content
	%\begin{frame}
		%\tableofcontents[sectionstyle=show,subsectionstyle=show/shaded/hide,subsubsectionstyle=show/shaded/hide]
	%\end{frame}
		


\begin{frame}{Question 2}
\vspace{0.5em}
\textbf{Problem 2(a):} Compute the $\ell$ such that:
\[
\frac{u_{\ell}(c,\ell)}{u_c(c,\ell)} - z f_h(k,h) = 0, \qquad \ell = 1-h.
\]

\vspace{0.5em}
\textbf{Problem 2(b):} Compute the $(h^*, s^*)$ such that:
\[
G(h,s) = 0
\]

$G(h,s)$ is the vector of intratemporal conditions for $h$ and $s$ derived from the household problem.

\vspace{0.5em}
Main idea: transfer economic problem into a root-finding problem and solve it robustly using numerical methods.
\end{frame}

\begin{frame}{General Solver Strategy}
I use the same routines from Part 1 by rewriting each problem as
\[
g(x)=0,
\]
Then use newton method to update $x$ iteratively:
\[
x^{i+1} = x^i - J_G(x^i)^{-1} G(x^i)
\]
since we don't have analytical derivatives for $g(\cdot)$, so I use numerical approximations:
\[
J_G(x) \approx \frac{G(x+\delta e_j) - G(x-\delta e_j)}{2\delta}.
\]
$e_j$ is the unit vector in the $j$-th direction.

\end{frame}

\begin{frame}{Problem 2(a): Economic System}
Define
\[
g(\ell)=\frac{u_{\ell}(c,\ell)}{u_c(c,\ell)}- z f_h(k,1-\ell).
\]
Then $\ell^*$ solves
\[
g(\ell^*)=0.
\]


\end{frame}


\begin{frame}{Pseudo-Code for Problem 2(a)}
\small

\textbf{Input:} $c$, $z$, $k$, $u(c,\ell)$, $f(k,h)$, initial guess $\ell^0$, tol, max\_iter, $\delta$.

\textbf{Output:} $\ell^*$, number of iterations, convergence status


Choose initial guess $\ell_0$. Check whether $\ell_0 \in [0,1]$

For iteration $n=0,1,2,\dots$:
\begin{enumerate}
    \item Evaluate $g(\ell_n)$.
    \item Approximate $g'(\ell_n)$ numerically.
    \item If $|g'(\ell_n)|$ is too small, stop with diagnostic message.
    \item Update
    \[
    \ell_{n+1}=\ell_n-\frac{g(\ell_n)}{g'(\ell_n)}.
    \]
    \item stop if $\ell_{n+1} \in [0,1]$, return infeasible warning
    \item Stop if
    \[
    |g(\ell_{n+1})|<\texttt{tol}.
    \]
\end{enumerate}
Return $\ell^*$ and $h^*=1-\ell^*$.

\end{frame}

\begin{frame}{Robustness Design in Problem 2(a)}

\begin{enumerate}
    \item \textbf{Feasibility check for leisure:}
    \[
    0 \le \ell \le 1.
    \]
    This ensures labor supply $h=1-\ell$ is also feasible.

    \item \textbf{Check that marginal utility is well-defined:}
    \[
    u_c(c,\ell) \neq 0 \implies \frac{u_{\ell}(c,\ell)}{u_c(c,\ell)} \quad \text{well defined}
    \]

    \item \textbf{Check $g(\ell)$ and $g'(\ell)$:}

    \begin{align*}
        \bigl(g(\ell_n)\in\mathbb{R}\bigr)\land\bigl(g'(\ell_n)\neq 0\bigr)\land\bigl(g'(\ell_n)\in\mathbb{R}\bigr)
        \;\implies\;
        \ell_{n+1}=\ell_n-\frac{g(\ell_n)}{g'(\ell_n)}.
    \end{align*}

    \item \textbf{Maximum-iteration rule:} if convergence is not reached within the iteration limit, the solver reports failure explicitly.
\end{enumerate}
\end{frame}

\begin{frame}{Problem 2(b): Economic System}
Let the unknowns be work hours $h$ and shopping hours $s$, with leisure
\[
\ell = 1-h-s.
\]

Write the equilibrium conditions as a two-dimensional system:
\[
G(h,s)=
\begin{bmatrix}
G_1(h,s)\\[0.3em]
G_2(h,s)
\end{bmatrix}
=
0.
\]
solve $(h^*,s^*)$ such that $G(h^*,s^*) = 0$

\end{frame}


\begin{frame}{Pseudo-Code for Problem 2(b)}
\small

\textbf{Input:} $w$, $y$, $u(c,\ell)$, $p(s)$, initial guess $(h^0,s^0)$, tol, max\_iter, $\delta$.

\textbf{Output:} $(h^*, s^*)$, number of iterations, convergence status

Choose initial guess
\[
x_0=
\begin{bmatrix}
h_0\\
s_0
\end{bmatrix}
\]

check initial feasibility:
\[
h^0 \ge 0,\qquad s^0 \ge 0,\qquad 1-h^0-s^0 \ge 0.
\]

For iteration $n=0,1,2,\dots$:
\begin{enumerate}
    \item Evaluate residual vector $G(x_n)$.
    \item Construct numerical Jacobian $J_G(x_n)$, then update
    \[
    x_{n+1}=x_n - J_G(x_n)^{-1} G(x_n).
    \]
    \item check feasibility, stop and return a warning if:
    \[
    h_{n+1} \ge 0,\qquad s_{n+1} \ge 0,\qquad 1-h_{n+1}-s_{n+1} \ge 0.
    \]
    \item Stop if
    \[
    \|G(x_{n+1})\|<\texttt{tol}.
    \]
\end{enumerate}
Return
\[
x^*=
\begin{bmatrix}
h^*\\
s^*
\end{bmatrix}.
\]

\end{frame}

\begin{frame}{Robustness Design in Problem 2(b)}


\begin{enumerate}
    \item \textbf{Feasible time allocation:}
    \[
    h \ge 0,\qquad s \ge 0,\qquad \ell = 1-h-s \ge 0.
    \]

    \item \textbf{Positive price:}
    \[
    p(s) > 0.
    \]

    \item \textbf{Check $G(h,s)$ and $J_G(h,s)$:}

    \begin{align*}
        \bigl(G(h_n,s_n)\in\mathbb{R}\bigr)\land\bigl(J_G(h_n,s_n)\in\mathbb{R}\bigr)
        \land\bigl(J_G(h_n,s_n)\neq 0\bigr) \land\bigl(J_G(x^i) \quad \text{singular}\bigr) \\
        \;\implies\;
        x_{n+1}=x_n - J_G(x_n)^{-1} G(x_n)
    \end{align*}

    \item \textbf{Maximum-iteration rule:} if convergence is not reached within the iteration limit, the solver reports failure explicitly.
\end{enumerate}
\end{frame}


\begin{frame}{Conclusion}
\begin{itemize}
    \item Problem 2(a) was solved as a \textbf{scalar} root-finding problem using numerical derivatives.
    \item Problem 2(b) was solved as a \textbf{vector} root-finding problem using a numerical Jacobian.
    \item The main robustness features are:
    \begin{itemize}
        \item feasibility checks
        \item finite-value checks
        \item derivative / Jacobian checks
        \item explicit convergence diagnostics
    \end{itemize}
\end{itemize}

\end{frame}
\end{document}